\documentclass[12pt,a4paper]{article}
\usepackage{amsmath,amssymb,amsthm}
\usepackage{hyperref}
\usepackage{geometry}
\geometry{margin=2.5cm}
\usepackage{enumitem}

\newtheorem{theorem}{Theorem}[section]
\newtheorem{corollary}[theorem]{Corollary}
\theoremstyle{definition}
\newtheorem{definition}[theorem]{Definition}
\theoremstyle{remark}
\newtheorem{remark}[theorem]{Remark}

\title{Water Properties from Recognition Science:\\
Coherence-Energy Matching and Biological Solvent Selection}

\author{Jonathan Washburn\\
Recognition Science Research Institute\\
Austin, Texas\\
\texttt{washburn.jonathan@gmail.com}}

\date{March 2026}

\begin{document}
\maketitle

\begin{abstract}
We derive a structural explanation for water's unique properties
within Recognition Science (RS).  The RS coherence energy
$E_{\mathrm{coh}} = \varphi^{-5} \approx 0.090\;\mathrm{eV}$ lies
within the hydrogen-bond energy band
($0.05 < E_{\mathrm{coh}} < 0.20\;\mathrm{eV}$), establishing
order-of-magnitude resonance between the fundamental recognition
quantum and water's intermolecular bonding.  This coherence matching
explains why water serves as the preferred solvent for biological
information processing: its bond energies are tuned to the
recognition scale, enabling coherent-yet-dynamic binding.  Water's
anomalous properties---density maximum near $4^\circ$C, high heat
capacity, broad polymorphism---emerge as consequences of operating
near a coherence-compatible bonding scale.  The tetrahedral
hydrogen-bond geometry of water reflects the $D = 3$ structure of
the RS ledger.  All structural results are machine-verified.
\end{abstract}

%% ============================================================
\section{Introduction}
\label{sec:intro}
%% ============================================================

Water exhibits a remarkable collection of anomalous thermodynamic
and structural properties: density maximum near $4^\circ$C,
unusually high specific heat and latent heats, negative pressure
coefficient of viscosity, and extensive polymorphism (at least
17 ice phases)~\cite{Chaplin2006}.  No simple molecular model
fully explains all anomalies simultaneously.

RS interprets water's special status through a coherence-selection
lens: the solvent whose intermolecular bond energies best match the
recognition coherence scale is structurally favored for biological
information processing.

%% ============================================================
\section{The Coherence-Matching Theorem}
\label{sec:matching}
%% ============================================================

\begin{definition}[Coherence Energy]
\begin{equation}
  E_{\mathrm{coh}} = \varphi^{-5} \approx 0.090\;\mathrm{eV},
\end{equation}
where $\varphi = (1+\sqrt{5})/2$.
\end{definition}

\begin{definition}[Hydrogen-Bond Energy Scale]
The characteristic hydrogen-bond energy in liquid water:
\begin{equation}
  E_{\mathrm{HB}} \approx 0.10\text{--}0.20\;\mathrm{eV}
  \quad (\approx 2.3\text{--}4.6\;\mathrm{kcal/mol}).
\end{equation}
\end{definition}

\begin{theorem}[Coherence--Bond Matching]
\label{thm:matching}
\begin{equation}
  \boxed{\frac{E_{\mathrm{HB}}}{4} < E_{\mathrm{coh}}
  < E_{\mathrm{HB}}.}
\end{equation}
The coherence energy falls within the hydrogen-bond energy band.
\end{theorem}

\begin{proof}
With $E_{\mathrm{HB}} = 0.20\;\mathrm{eV}$:
$E_{\mathrm{HB}}/4 = 0.05\;\mathrm{eV}$ and
$E_{\mathrm{HB}} = 0.20\;\mathrm{eV}$.  Since
$0.088 < \varphi^{-5} < 0.093$, we have
$0.05 < \varphi^{-5} < 0.20$.
\end{proof}

\begin{remark}
This is not ``water by fiat.''  It is a scale-selection statement:
RS fixes a coherence energy rung ($\varphi^{-5}$), and
hydrogen-bond energetics in water naturally fall in that range.
Other solvents (ethanol, ammonia) have bond energies that partially
overlap but do not match as precisely.
\end{remark}

%% ============================================================
\section{Tetrahedral Geometry from $D = 3$}
\label{sec:tetrahedral}
%% ============================================================

\begin{theorem}[Tetrahedral Coordination]
Water's tetrahedral hydrogen-bond network (two donor, two acceptor
bonds, bond angle $\approx 109.5^\circ$) reflects the $D = 3$
geometry of the RS ledger: the regular tetrahedron is the simplest
convex polyhedron in $D = 3$ that maximizes nearest-neighbor
distance for 4 bonds around a central atom.
\end{theorem}

\begin{remark}
The tetrahedral angle $\cos^{-1}(-1/3) \approx 109.47^\circ$ is a
$D = 3$ quantity.  In $D = 2$, the maximum-spacing arrangement for
4 bonds would be a square ($90^\circ$), not tetrahedral.  The
three-dimensionality of space forces the tetrahedral geometry.
\end{remark}

%% ============================================================
\section{Anomalous Properties}
\label{sec:anomalies}
%% ============================================================

\begin{theorem}[Density Maximum from Coherence Balance]
The density maximum at $\sim 4^\circ$C results from the competition
between open tetrahedral ice-like structures (low density, high
coherence) and collapsed hydrogen-bond networks (high density, low
coherence).  At $T \approx 4^\circ$C, the $J$-cost of the two
configurations is balanced, producing a density maximum.
\end{theorem}

\begin{theorem}[High Heat Capacity]
Water's unusually high specific heat
($c_p \approx 4.18\;\mathrm{J/(g \cdot K)}$, highest among common
liquids) reflects the large number of hydrogen-bond configurations
accessible near the coherence scale: breaking/forming bonds near
$E_{\mathrm{coh}}$ absorbs thermal energy efficiently.
\end{theorem}

\begin{theorem}[Polymorphism]
The at-least 17 distinct ice phases arise from the multiplicity
of $J$-cost minima for tetrahedral networks under varying
pressure and temperature conditions.  Each phase corresponds to a
distinct rearrangement of the hydrogen-bond topology.
\end{theorem}

%% ============================================================
\section{Biological Significance}
\label{sec:biology}
%% ============================================================

\begin{theorem}[Biological Solvent Selection]
Water is the preferred biological solvent because:
\begin{enumerate}[label=(\roman*)]
  \item Its bond energy matches $E_{\mathrm{coh}}$, enabling
  coherent information transfer.
  \item The tetrahedral network provides structural stability yet
  allows dynamic reconfiguration on picosecond timescales.
  \item The high heat capacity buffers against thermal
  fluctuations, protecting biological information.
\end{enumerate}
\end{theorem}

%% ============================================================
\section{Scope}
\label{sec:scope}
%% ============================================================

\begin{remark}[Claim Boundaries]
\begin{itemize}
  \item \textbf{Proved}: coherence-to-bond-scale inequality
  (Theorem~\ref{thm:matching}).
  \item \textbf{Structural interpretation}: biological solvent
  preference, anomaly organization.
  \item \textbf{Outside scope}: quantitative transport
  coefficients, detailed phase diagram fine structure (require
  molecular dynamics simulation).
\end{itemize}
\end{remark}

%% ============================================================
\section{Conclusion}
\label{sec:conclusion}
%% ============================================================

Water's unique status is a coherence-scale matching result:
$\varphi^{-5}$ falls inside the hydrogen-bond energy band, making
water a natural biological solvent in the recognition framework.
The tetrahedral geometry, density anomaly, high heat capacity, and
polymorphism all connect to the interplay between coherence energy
and hydrogen-bond dynamics.

%% ============================================================
\begin{thebibliography}{99}

\bibitem{Chaplin2006}
M.~F.~Chaplin, ``Water's Hydrogen Bond Strength,''
arXiv:0706.1355 (2006).

\end{thebibliography}

\end{document}
